% !PRINTED 157
\weberpage{157}

This expression now shows that $R$, apart from the factor
$(-1)^{mn}$, depends on $\varphi(x)$ in the same way as
on $f(x)$, and that we may therefore also set

\weq{7}{(-1)^{mn}\, R = b_0^n\, f(\beta_1)\, f(\beta_2)\, \ldots\, f(\beta_m).}

The discriminants are accordingly a special case of the
resultants, which one obtains when one takes for $\varphi(x)$ the derivative
$f'(x)$.

\gesperrt{The vanishing of the resultant of $f(x)$ and
$\varphi(x)$ is the necessary and sufficient condition
for the two equations}

\[ f(x) = 0 \quad \text{and} \quad \varphi(x) = 0 \]

\gesperrt{to have a common root.}

The equation that arises by setting the resultant equal to zero
is also called the \gesperrt{final equation} from $f(x) = 0$ and $\varphi(x) = 0$,
and the setting up of the \gesperrt{final equation} the \gesperrt{elimination}
of $x$.

For the computation of resultants one may apply the algorithm
of the greatest common divisor (\S.\,6). For if
$m \leqq n$, one can determine an integral function $Q$ and a like
 function $\psi(x)$, whose degree is smaller than $m$, in such a way
 that

\[ f(x) = Q\varphi(x) + \psi(x), \]

and hence, if one puts for $x$ one of the roots $\beta_i$,

\[ f(\beta_i) = \psi(\beta_i), \]

hence

\weq{8}{b_0^n \overset{i}{\Pi} f(\beta_i) = b_0^n \overset{i}{\Pi} \psi(\beta_i);}

if $m'$ is the degree of $\psi(x)$, then

\weq{9}{b_0^{m'} \overset{i}{\Pi} \psi(\beta_i) = R'}

is the resultant of $\varphi(x)$ and $\psi(x)$, and therefore

\weq{10}{R = b_0^{n-m'}\, R'.}

The formation of the resultant $R$ is hereby reduced to the formation
of $R'$. One can now divide $\varphi(x)$ again by $\psi(x)$,
and thus reduces the formation of $R'$ to the formation of a
resultant of functions of ever lower degree, until
one finally arrives at a constant, i.\,e. a function of the
coefficients $a$, $b$, which must be identical with the resultant
 sought.
